%-*-latex-*-
%-*-latex-*-
\newcommand\COURSE{ciss245}
\newcommand\ASSESSMENT{q6101}
\newcommand\ASSESSMENTTYPE{Quiz}
\newcommand\POINTS{\textwhite{xxx/xxx}}

\input{myquizpreamble}
\input{yliow}
\input{\COURSE}
%-*-latex-*-

%-*-latex-*-

%-*-latex-*-
\renewcommand\TITLE{\ASSESSMENTTYPE \ \ASSESSMENT}

\renewcommand\EMAIL{}
\renewcommand\AUTHOR{YOUR EMAIL}

\textwidth=6in
\begin{document}
\topmatterthree

%------------------------------------------------------------------------------
\nextq
Wumpus lives in the Wumpus world, which is a 4-by-4 grid.
Initially, Wumpus is at any cell of the grid except for (3, 0) (row 3,
column 0).
A hunter also lives in the Wumpus world.
He/she walks in the direction he/she is facing.
Initially he/she is facing N (north).
In each time step in Wumpus world, the hunter
(1) randomly selects a direction,
(2) randomly decides to walk one step (in his/her direction) or not,
and
(3) randomly decides to fire one arrow (in his/her direction) or not.
Here's a test run:
\begin{Verbatim}[frame=single,fontsize=\footnotesize]
+-+-+-+-+
|W| | | |
+-+-+-+-+
| | | | |
+-+-+-+-+
| | | | |
+-+-+-+-+
|H| | | |
+-+-+-+-+
Hunter direction: N

+-+-+-+-+
| | | | |
+-+-+-+-+
|W| | | |
+-+-+-+-+
| | | | |
+-+-+-+-+
| |H| | |
+-+-+-+-+
Hunter direction: E

+-+-+-+-+
| | | | |
+-+-+-+-+
| |W| | |
+-+-+-+-+
| | | | |
+-+-+-+-+
| |H| | |
+-+-+-+-+
Hunter direction: S

+-+-+-+-+
| |W| | |
+-+-+-+-+
| |^| | |
+-+-+-+-+
| |H| | |
+-+-+-+-+
| | | | |
+-+-+-+-+
Hunter direction: N

+-+-+-+-+
| |w| | |
+-+-+-+-+
| | | | |
+-+-+-+-+
| | |H| |
+-+-+-+-+
| | | | |
+-+-+-+-+
Hunter direction: E

+-+-+-+-+
| |w| | |
+-+-+-+-+
| | | | |
+-+-+-+-+
| | | |H|
+-+-+-+-+
| | | | |
+-+-+-+-+
Hunter direction: E
\end{Verbatim}
The \verb!^! is the arrow flying in the N direction.
You should draw the arrow using
\verb!^!,
\verb!v!,
\verb!>!,
\verb!<! depending on the direction
N,S,E,W (respectively)
of flight.

Note that when the hunter fires his/her arrow, the arrow will fly in the
direction of the hunter (at that point in time).
The hunter can only fire one arrow (too bad).

When the arrow kills Wumpus (i.e., the arrow and Wumpus are
at the same cell of the grid), a \verb!w! (lowercase of W)
is printed at that cell.
After that, Wumpus is dead and does not move (i.e., \verb!w! does not
move).
If hunter is in the same cell as the dead Wumpus, \verb!H! is printed.

If Wumpus and the hunter are both in the same cell,
Wumpus kills the hunter (sad),
\verb!W! is printed at that cell.
When Wumpus moves away from that cell, \verb!h! (the dead hunter)
is printed.

You are given this:
\begin{Verbatim}[frame=single,fontsize=\footnotesize]
// file: WumpusWorld.h
#ifndef WUMPUSWORLD_H
#define WUMPUSWORLD_H

#include <iostream>

class WumpusWorld
{
public:
    WumpusWorld();       // initialize so that Wumpus is at any cell other than
                         // (3,0).
                         // hunter is at (3,0).
    void println();      // print according to the above format
    void move_wumpus();  // randomize a direction rand() % 4 where 0,1,2,3
                         // are N,S,E,W for Wumpus. Of course Wumpus
                         // must stay in the world.
                         // If wumpus tries to move N but its row is 0,
                         // wumpus stays put.
                         // A dead Wumpus does not move.
    void move_hunter();  // 1. randomize a direction rand() % 4 where 0,1,2,3
                         //    are N,S,E,W for hunter.
                         // 2. randomize whether to walk or not
                         // 3. randomize whether to fire arrow or not
                         // Of course hunter must stay in the world.
                         // If hunter tries to move N but its row is 0,
                         // hunter stays put.
                         // A dead hunter does not move.
    void move_arrow();   // move arrow if necessary
private:
    char world_[4][4];
    int hunter_direction_;
    int arrow_state_;    // 0-N, 1-S, 2-E, 3-W, 4-not fired, 5-hit something
};
#endif
\end{Verbatim}

You want to make the hunter wonder around like this:
\begin{Verbatim}[frame=single,fontsize=\footnotesize]
#include <iostream>
#include <ctime>
#include <cstdlib>
#include "WumpusWorld.h"

int main()
{
    srand((unsigned int) time(NULL));

    WumpusWorld ww;
    ww.println();
    for (int i = 0; i < 50; ++i)
    {
        ww.move_wumpus();
        ww.move_hunter();
        ww.move_arrow();
        ww.println();
    }
    return 0;
}
\end{Verbatim}
Complete the following cpp file that contains the implementation
of the methods declared in the WumpusWorld class.
\\
\ANSWER
\begin{answercode}
// file: WumpusWorld.cpp

\end{answercode}

%------------------------------------------------------------------------------
%-*-latex-*-
\newpage
\input{instructions}
\end{document}
